% ── MAIN MATTER ─────────────────────────────────────────────
\mainmatter
\pagestyle{fancy}

% ============================================================
%  CHAPTER 1: INTRODUCTION
% ============================================================
\chapter{مقدمه}
\label{ch:intro}

ایران در مقطعی تاریخی قرار دارد که تصمیمات امروز، سرنوشت نسل‌های آینده را رقم خواهد زد. در چنین لحظه‌ای، نه تعلل و بی‌عملی کارساز است و نه شتاب‌زدگی، هیجان و تصمیم‌های احساسی. همچنین تلاش برای مصادره‌ی جنبش یا هدایت آن به سود فرد یا گروهی خاص، می‌تواند آینده‌ی کشور را با مخاطره روبه‌رو کند.
برای دستیابی به ایرانی آزاد، آباد، متکثر و چندصدایی — که پایه‌های مستحکم انسجام اجتماعی، آینده‌نگری و توسعه‌ی پایدار در آن استوار باشد — لازم است فرآیند گذار با دقت، مسئولیت‌پذیری و با بهره‌گیری از درس‌های آموخته‌شده‌ی ملی و بین‌المللی طراحی و هدایت شود.

خشم اجتماعی و میل به اقدام سریع برای عبور از جمهوری اسلامی، تا حدی قابل درک و حتی در برخی شرایط اجتناب‌ناپذیر است؛ به‌ویژه هنگامی که معطوف به اطمینان از بازگشت‌ناپذیری رژیم سفاک جمهوری اسلامی باشد. با این حال، خشونت کور یا اقدامات خارج از چارچوب، نه‌تنها راه‌گشا نیست، بلکه می‌تواند مانع شکل‌گیری همبستگی ملی و همکاری میان همه‌‌ی کسانی شود که از این نظام آسیب دیده‌اند و برای آینده‌ای بهتر تلاش می‌کنند.


\begin{summarybox}
تجربه‌ی جهانی — از تونس و آفریقای جنوبی تا اسپانیا، لهستان و شیلی، و نیز تجربه‌های ناکام در کشورهایی چون لیبی، مصر و عراق — نشان می‌دهد که 	\textbf{کیفیت طراحی نهادهای دوران گذار}  بیش از هر عامل دیگری در تعیین سرنوشت نظام سیاسی آینده نقش دارد. کشورهایی که این مرحله را با آگاهی، مشارکت و اجماع ملی مدیریت کرده‌اند، توانسته‌اند از دام هرج‌ومرج، بازگشت استبداد و حتی جنگ داخلی پرهیز کنند.
\end{summarybox}

% ============================================================
%  CHAPTER 2: GUIDING PRINCIPLES
% ============================================================
\chapter{اصول راهنمای طراحی نظام نهادی گذار}
\label{ch:principles}

پیش از ورود به جزئیات ساختاری، لازم است اصول بنیادینی که طراحی کل نظام نهادی گذار
بر مبنای آن‌ها صورت گرفته، به‌روشنی بیان شود. این اصول از ادبیات نظری سیاست
تأسیسی (\lr{Constituent Politics})، تجربه‌های عملی گذار دموکراتیک در جهان 
نیز با توجه به شرایط و ضرورت‌های خاص ایران استخراج شده‌اند. این اصول نه تنها چارچوبی برای طراحی نهادها و فرایندهای گذار فراهم می‌کنند، بلکه معیارهایی برای ارزیابی تصمیمات و اقدامات در طول این دوره‌ی حساس هستند. در ادامه، هر اصل به‌طور جداگانه بیان و به شکل خلاصه‌ای با نمونه‌های جهانی مرتبط توضیح داده شده است.\par
سیاست تأسیسی یا (\lr{Constituent Politics}), در واقع به سیاست‌های پایه‌ای و بنیادینی اشاره دارد که ساختار، قوانین، و اصول بنیادین یک نهاد سیاسی (مثل دولت) را تشکیل می‌دهند. این نوع سیاست، برعکس سیاست‌های عادی (قانون‌گذاری معمولی)، قوانین بازی سیاسی را برای آینده تعیین می‌کند، از جمله تدوین قانون اساسی، ساختار نهادها، و حقوق اساسی شهروندان. در دوران گذار دموکراتیک، سیاست تأسیسی اهمیت ویژه‌ای پیدا می‌کند، زیرا تصمیمات گرفته شده در این مرحله، مسیر توسعه‌ی سیاسی، اجتماعی و اقتصادی کشور را برای دهه‌ها یا حتی قرن‌ها تعیین می‌کند. اصول راهنمای طراحی نظام نهادی گذار، در واقع چارچوبی برای سیاست تأسیسی در ایران فراهم می‌کنند تا اطمینان حاصل شود که این فرایند به شکلی عادلانه، فراگیر و پایدار انجام شود.


% --- Principle 1 ---
\needspace{6\baselineskip}
\begin{principlebox}{اصل اول: ترکیب دموکراسی مستقیم و تخصص‌گرایی}
	\begin{quotebox}
		هر تصمیم سیاسی دارای دو بُعد اساسی است: بُعد واقعیت (\lr{Fact}) و بُعد ارزش (\lr{Value}).
		هنر حکمرانی خوب بخصوص در دوران گذار، که مقطع حساس‌تری است، جدا کردن و سپس ترکیب درست این دو بُعد است.
	\end{quotebox}
	
	\vspace{0.5em}
	\begin{itemize}
		\item \textbf{بُعد واقعیت و تخصص (\lr{What is})}: داده‌ها، حقایق علمی‌ - تجربی،
		تحلیل‌های فنی، پیامدهای اقتصادی، حقوقی و اجتماعی. این بُعد از \textbf{تکنوکرات‌ها،
			متخصصان و مراکز پژوهشی} دریافت می‌شود.
		
		\item \textbf{بُعد ارزش و خواست ملی (\lr{What ought to be})}: تمایلات، ارزش‌ها،
		آرزوها، ترجیحات و حق تعیین سرنوشت مردم. این بُعد \textbf{فقط و فقط} از طریق
		\textbf{مراجعه‌ی مستقیم به آرای مردم} قابل تعیین است.
	\end{itemize}
	
	\vspace{0.5em}
	\textbf{مبنای نظری:} این تفکیک ریشه در اندیشه‌ی دیوید هیوم (تفکیک «هست» از «باید»)،
	یورگن هابرماس (دموکراسی مشورتی)، و فیلیپ پتیت (جمهوری‌خواهی و آزادی به‌مثابه
	عدم سلطه) دارد.
	
	\textbf{نمونه‌های جهانی:}
	\begin{itemize}
		\item \textbf{ایرلند (۲۰۱۶--۲۰۱۹)}: مجمع شهروندان (\lr{Citizens' Assembly}) --- ترکیب شهروندان عادی و متخصصان؛ نتیجه: اصلاح قانون اساسی از جمله قانون سقط جنین
		\item \textbf{فرانسه (۲۰۱۹--۲۰۲۰)}: کنوانسیون شهروندی اقلیم --- ۱۵۰ شهروند تصادفی با کمک متخصصان اقلیمی
		\item \textbf{آفریقای جنوبی}: تدوین قانون اساسی ۱۹۹۶ شامل بیش از ۲ میلیون نظر مردمی
	\end{itemize}
\end{principlebox}

% --- Principle 2 ---
\needspace{6\baselineskip}
\begin{principlebox}{اصل دوم: نمایندگی نسبی (\lr{Proportional Representation})}
	\begin{quotebox}
		در جامعه‌ای به تنوع و تکثر ایران، هیچ صدایی نباید نشنیده بماند.
	\end{quotebox}
	
	سیستم انتخاباتی نسبی‌فهرستی تضمین می‌کند که ترکیب مجلس مهستان آینه‌ی واقعی جامعه‌ی
	ایران باشد. تجربه‌ی جهانی (لیپهارت، ۱۹۹۹؛ رینولدز، ۲۰۱۱) نشان داده که سیستم‌های
	نسبی در جوامع متکثر:
	\begin{itemize}
		\item نمایندگی عادلانه‌تر اقلیت‌ها را تضمین می‌کنند
		\item ائتلاف‌سازی و مذاکره را تشویق می‌کنند
		\item ریسک انحصار قدرت توسط یک گروه را کاهش می‌دهند
		\item مشارکت سیاسی را افزایش می‌دهند
	\end{itemize}
\end{principlebox}

% --- Principle 3 ---
\needspace{6\baselineskip}
\begin{principlebox}{اصل سوم: قانون اساسی‌گرایی دو مرحله‌ای}
	\begin{quotebox}
		خلأ قانونی خطرناک‌ترین تهدید دوران گذار است.
	\end{quotebox}
	
	مدل دو قانون اساسی (موقت + نهایی) --- بر اساس تجربه‌ی موفق آفریقای جنوبی ---
	تضمین می‌کند که:
	\begin{itemize}
		\item از روز اول چارچوب حقوقی وجود داشته باشد
		\item فشار زمانی برای تصمیمات سرنوشت‌ساز حذف شود
		\item اصول غیرقابل‌نقض از ابتدا تثبیت شوند
		\item فرصت کافی برای بحث عمومی وجود داشته باشد
	\end{itemize}
\end{principlebox}

% --- Principle 4 ---
\needspace{6\baselineskip}
\begin{principlebox}{اصل چهارم: اصلاح نه تخریب}
	\begin{quotebox}
		تجربه‌ی عراق نشان داد: انحلال کامل نهادها = فاجعه.
	\end{quotebox}
	
	ساختارهای اداری، نظامی و خدماتی کشور با \textbf{رهبری جدید و اصلاح‌شده} حفظ می‌شوند.
	نیروهای حرفه‌ای و غیرمتخلف جذب ساختار جدید شده و فقط ساختارهای ایدئولوژیک و عاملان
	جنایات کنار گذاشته می‌شوند. سیاست \lr{De-Baathification} عراق درس آموخته‌ای است
	که این طرح از آن عبرت گرفته است.
\end{principlebox}

% --- Principle 5 ---
\needspace{6\baselineskip}
\begin{principlebox}{اصل پنجم: شفافیت و پاسخ‌گویی حداکثری}
	\begin{quotebox}
		قدرت بدون نظارت، فاسد می‌شود.
	\end{quotebox}
	
	تمام جلسات مجلس مهستان علنی خواهد بود. تمام تصمیمات مستند و عمومی خواهند شد.
	تمام مقامات دوران گذار ملزم به شفافیت مالی هستند.
\end{principlebox}

% --- Principle 6 ---
\needspace{6\baselineskip}
\begin{principlebox}{اصل ششم: فراگیری و تعهد به حقوق بنیادین}
	\begin{quotebox}
		ایران برای همه‌ی ایرانیان.
	\end{quotebox}
	
	هیچ گروه، قوم، جنسیت، مذهب یا ایدئولوژی از فرایند گذار حذف نخواهد شد. اصول بنیادین
	حقوق بشر --- شامل برابری جنسیتی، حقوق اقلیت‌ها، آزادی عقیده و ممنوعیت شکنجه ---
	\textbf{غیرقابل‌نقض} هستند و هیچ نهادی، حتی مجلس مهستان، حق تعلیق آن‌ها را ندارد.
\end{principlebox}

% ============================================================
%  CHAPTER 3: MAHESTAN ASSEMBLY FRAMEWORK
% ============================================================
\chapter{چارچوب گذار: مجلس مهستان}
\label{ch:mahestan}

\begin{notebox}[title={\rl{یادداشت مهم درباره‌ی اعداد و نسبت‌های این سند}}]
	تمام اعداد، نسبت‌ها، آستانه‌ها و زمان‌بندی‌های مشخصی که در این سند ذکر می‌شوند
	(مانند ۳۰۰ نماینده، ۳\% آستانه، ۲ ماه مهلت) \textbf{پیشنهادی برای بحث} هستند ---
	نه مقادیر ثابت. منظور اصلی این سند، نشان دادن \emph{نیاز به تعریف چنین پارامترهایی}
	است. تمام ارقام توسط مجلس مهستان منتخب مردم نهایی خواهند شد.
\end{notebox}

\textbf{نهاد محوری گذار} یک \textbf{مجلس مردم‌محور منتخب} (مجلس مهستان) است که
در اسرع وقت پس از آغاز دوره‌ی گذار تشکیل خواهد شد.

% ──────────────────────────────────────────────────────────
\section{فرایند انتخابات و تشکیل مجلس}
% ──────────────────────────────────────────────────────────

\begin{bandbox}{بند ۱}
	جریان‌ها و ائتلاف‌هایی به شکل خودجوش باید شکل گرفته و اعضایی را در کل کشور به شکل
	فهرست کاندیداهای مورد تأیید معرفی کنند. \textbf{هر جریانی که فهرست ملی معرفی می‌کند،
		موظف است پیش‌نویس قانون اساسی موقت پیشنهادی خود را نیز منتشر کند.}
\end{bandbox}

\begin{bandbox}{بند ۲}
	پس از آغاز دوره‌ی گذار، فرایند ثبت‌نام و تبلیغات رسمی‌تر آغاز می‌شود. تبلیغات و معرفی
	و ائتلاف‌ها از همین اکنون نیز قابل انجام است. ثبت‌نام رسمی افراد و معرفی فهرست‌ها در
	وب‌سایت وزارت کشور و سایر پایگاه‌های اطلاع‌رسانی عمومی منعکس خواهد شد.
\end{bandbox}

\begin{bandbox}{بند ۳}
	اجرای انتخابات به شکل \textbf{غیرمتمرکز} توسط فرمانداری‌ها و شهرداری‌ها با کمک
	مدارس و آموزش‌وپرورش صورت خواهد گرفت. نظارت از طریق \textbf{ناظران جریان‌ها + ناظران
		بین‌المللی + ناظران جامعه‌ی مدنی} تضمین می‌شود.
\end{bandbox}

\begin{bandbox}{بند ۴}
	انتخابات در زمان مقتضی صورت خواهد گرفت و پس از شمارش آرا و تأیید نتایج، ظرف یک هفته
	\textbf{مجلس مهستان دوران گذار} تشکیل خواهد شد.
\end{bandbox}

% ──────────────────────────────────────────────────────────
\section{ترکیب مجلس مهستان}
% ──────────────────────────────────────────────────────────

\begin{bandbox}{بند ۵}
	مجلس مهستان متشکل از \textbf{۳۰۰ نماینده} (پیشنهادی) خواهد بود:
	\begin{itemize}
		\item \textbf{۲۴۰ کرسی استانی}: بر اساس جمعیت هر استان، سیستم نسبی‌فهرستی
		(حداقل ۳ کرسی برای هر استان). آستانه‌ی ورود: ۳\%
		\item \textbf{۴۰ کرسی فهرست ملی}: برای جریان‌هایی که فهرست سراسری دارند.
		آستانه‌ی ورود: ۵\%
		\item \textbf{۲۰ کرسی تضمینی اقلیت‌ها}: برای اقلیت‌های قومی، زبانی و مذهبی.
		انتخاب توسط خود اقلیت‌ها
	\end{itemize}
\end{bandbox}

\begin{notebox}[title={\rl{یادداشت درباره‌ی سیستم انتخاباتی (\lr{MMP} در برابر فهرستی نسبی)}}]
	طرح حاضر برای مجلس مهستان از \textbf{سیستم فهرستی نسبی} استفاده می‌کند --- اجرای
	آن در شرایط گذار ساده‌تر و سریع‌تر است. سیستم \textbf{\lr{MMP}} (مختلط
	موازی‌جبرانی --- الگوی آلمانی)، که هم نمایندگی محلی و هم عدالت نسبی را فراهم
	می‌کند، برای قانون اساسی \textbf{نهایی} به‌عنوان گزینه‌ی برتر پیشنهاد می‌شود.
	\lr{MMP} نیازمند زیرساخت فنی و آموزش رأی‌دهندگان بیشتری است که در هفته‌های اول
	گذار فراهم نیست؛ اما در نظامی تثبیت‌شده، مزایای آشکاری دارد.
\end{notebox}

\begin{bandbox}{بند ۶ --- سهمیه‌ی جنسیتی: قاعده‌ی ۱ به ۳}
	رعایت \textbf{حداقل ۳۰\% تنوع جنسیتی} در فهرست‌های انتخاباتی الزامی است.
	
	\textbf{ضمانت اجرایی:} در میان هر سه نفر متوالی در فهرست، حداقل یک نفر باید از
	جنسیت کمتر‌نمایندگی‌شده باشد. اگر سه نفر اول همگی مرد باشند، \textbf{نفر سوم}
	با بالاترین رأیی که به یک زن داده شده جایگزین می‌شود.
	
	\textbf{فهرست‌هایی که از تنظیم خود طبق این قاعده سرباز زنند، از ثبت‌نام رسمی
		محروم خواهند شد.}
\end{bandbox}

\begin{bandbox}{بند ۶-الف --- حق رأی دیاسپورا}
	میلیون‌ها ایرانی خارج از کشور از حقوق و آرزوهای یکسان با هموطنان داخل برخوردارند.
	\textbf{پیشنهاد پایه} (باز است و توسط مجلس نهایی خواهد شد): ایرانیان خارج از کشور
	می‌توانند \textbf{به فهرست منتخب خود در یک حوزه‌ی انتخاباتی داخلی} رأی دهند.
	گزینه‌های جایگزین (کرسی‌های اختصاصی دیاسپورا، رأی‌گیری آنلاین مستقل) نیز قابل بررسی هستند.
\end{bandbox}

% ──────────────────────────────────────────────────────────
\section{هیأت رئیسه و ساختار داخلی}
% ──────────────────────────────────────────────────────────

\begin{bandbox}{بند ۷}
	اگر جریانی با فهرست واحد سراسری به \textbf{اکثریت مطلق} (۵۰\%+۱) برسد، ریاست
	جلسه‌ی اول را بر عهده خواهد داشت. در غیر این صورت، \textbf{ریاست سنی}
	تعیین‌کننده‌ی ریاست موقت مجلس خواهد بود.
\end{bandbox}

\begin{bandbox}{بند ۸ --- کمیسیون‌های تخصصی دائمی}
	مجلس مهستان بلافاصله کمیسیون‌های تخصصی زیر را تشکیل می‌دهد:
	\begin{itemize}
		\item کمیسیون قانون اساسی موقت
		\item کمیسیون عدالت انتقالی و حقوق بشر
		\item کمیسیون اصلاح بخش امنیتی و نظامی (\lr{SSR})
		\item کمیسیون اصلاح قضایی و حقوقی
		\item کمیسیون اقتصاد و بودجه
		\item کمیسیون آموزش، علم و فرهنگ
		\item کمیسیون سیاست خارجی و روابط بین‌الملل
		\item کمیسیون حقوق اقلیت‌ها و تنوع قومی‌زبانی
		\item کمیسیون محیط زیست و منابع طبیعی
		\item کمیسیون حقوق زنان و برابری جنسیتی
		\item کمیسیون رسانه و حقوق دیجیتال
	\end{itemize}
\end{bandbox}

\begin{center}
\begin{tikzpicture}[
    comm/.style={draw=border-gray, fill=bg-gray, rounded corners=3pt,
        minimum width=3.3cm, minimum height=0.7cm, align=center, font=\tiny\bfseries}
]
    \node[comm, fill=MahestanBlue!10]  at (0,1.2)   {\rl{قانون اساسی موقت}};
    \node[comm, fill=mahestan-red!10]   at (3.8,1.2)  {\rl{عدالت انتقالی و حقوق بشر}};
    \node[comm, fill=mahestan-orange!10]at (7.6,1.2)  {\rl{اصلاح امنیتی و نظامی}};
    \node[comm, fill=mahestan-purple!10]at (11.4,1.2) {\rl{اصلاح قضایی و حقوقی}};
    \node[comm, fill=mahestan-green!10] at (0,0.3)    {\rl{اقتصاد و بودجه}};
    \node[comm, fill=mahestan-teal!10]  at (3.8,0.3)  {\rl{آموزش، علم و فرهنگ}};
    \node[comm, fill=MahestanBlue!10]  at (7.6,0.3)  {\rl{سیاست خارجی}};
    \node[comm, fill=mahestan-gold!15]  at (11.4,0.3) {\rl{حقوق اقلیتها}};
    \node[comm, fill=mahestan-green!10] at (1.9,-0.6) {\rl{محیط زیست}};
    \node[comm, fill=mahestan-red!10]   at (5.7,-0.6) {\rl{حقوق زنان و برابری}};
    \node[comm, fill=mahestan-purple!10]at (9.5,-0.6){ \rl{رسانه و حقوق دیجیتال}};
\end{tikzpicture}
\end{center}

% ──────────────────────────────────────────────────────────
\section{تشکیل کابینه‌ی اجرایی موقت}
% ──────────────────────────────────────────────────────────

\begin{bandbox}{بند ۹}
	در صورت وجود جریان اکثریت، این جریان مسئول تشکیل \textbf{کابینه‌ی اجرایی موقت}
	خواهد شد.
\end{bandbox}

\begin{bandbox}{بند ۱۰}
	در صورت عدم وجود جریان اکثریت، \textbf{ائتلافی} از جریان‌ها یا افراد مستقل که
	به اکثریت قاطع می‌رسند، کابینه را تشکیل می‌دهند.
\end{bandbox}

\begin{bandbox}{بند ۱۱}
	در صورت عدم تشکیل هیچ ائتلافی، مجلس مهستان برای هر پست کابینه به‌طور جداگانه
	رأی‌گیری می‌کند و \textbf{افرادی که بالاترین رأی در هر سمت کسب کنند، به‌طور
		خودکار وزیر آن حوزه می‌شوند.}
	
	پس از تشکیل این کابینه، \textbf{رئیس کابینه} می‌تواند تغییر هر عضو را به مجلس
	پیشنهاد دهد، مشروط به کسب \textbf{رأی اعتماد اکثریت ساده.}
\end{bandbox}

\begin{bandbox}{بند ۱۲ --- پاسخ‌گویی کابینه}
	کابینه به مجلس مهستان \textbf{پاسخ‌گو} است:
	\begin{itemize}
		\item رأی اعتماد اولیه: \textbf{۵۰\%+۱} نمایندگان
		\item استیضاح و برکناری: \textbf{۷۵\%} آرای نمایندگان
		\item تعویض فردی وزرا: \textbf{\lr{$\frac{2}{3}$}} آرا
	\end{itemize}
\end{bandbox}

\begin{center}
\begin{tikzpicture}[
    node distance=0.6cm,
    inst/.style={draw=#1, fill=#1!8, rounded corners=4pt,
        minimum width=4cm, minimum height=0.8cm, align=center, font=\small}
]
    % مرکز
    \node[draw=mahestan-green, fill=bg-light-green, rounded corners=6pt,
        minimum width=6cm, minimum height=1.2cm, align=center,
        font=\bfseries] (meh) {\rl{مجلس مهستان}\\[-2pt]{\footnotesize \rl{نهاد عالی حاکمیت دوران گذار}}};
    % ردیف اول
    \node[inst=mahestan-gold, below left=0.8cm and 0.5cm of meh]  (cab)
        {\rl{کابینه‌ی اجرایی موقت}};
    \node[inst=mahestan-red, below=0.8cm of meh]  (jud)
        {\rl{شورای قضایی انتقالی}};
    \node[inst=mahestan-orange, below right=0.8cm and 0.5cm of meh]  (mil)
        {\rl{فرماندهی نظامی-انتظامی}};
    % ردیف دوم
    \node[inst=mahestan-purple, below=0.5cm of cab]  (ssr) {\rl{کمیسیون اصلاح امنیتی}};
    \node[inst=mahestan-blue, below=0.5cm of jud]  (trc) {\rl{کمیسیون حقیقت‌یابی}};
    \node[inst=mahestan-teal, below=0.5cm of mil]  (elec) {\rl{هیأت نظارت انتخابات}};
    % ردیف سوم
    \node[inst=mahestan-green, below=0.5cm of ssr]  (omb) {\rl{آمبودزمان}};
    \node[inst=mahestan-blue, below=0.5cm of trc]  (med) {\rl{کمیسیون ملی رسانه‌ها}};
    \node[inst=mahestan-orange, below=0.5cm of elec]  (econ) {\rl{شورای اقتصادی}};
    % فلشها
    \draw[-Stealth, thick, mahestan-green] (meh) -- (cab);
    \draw[-Stealth, thick, mahestan-green] (meh) -- (jud);
    \draw[-Stealth, thick, mahestan-green] (meh) -- (mil);
\end{tikzpicture}
\end{center}
% ──────────────────────────────────────────────────────────
\section{مأموریت‌های ویژه و نهادسازی}
% ──────────────────────────────────────────────────────────

\subsection{الف) قانون اساسی موقت --- مأموریت فوری}

\begin{bandbox}{بند ۱۴}
	مجلس مهستان موظف است ظرف حداکثر \textbf{۲ ماه} پس از تشکیل:
	\begin{enumerate}
		\item \textbf{کمیسیون ویژه‌ی قانون اساسی موقت} را تشکیل دهد (۲۱ عضو مجلس + ۷ حقوقدان مستقل + ۳ ناظر بین‌المللی)
		\item پیش‌نویس‌های جریان‌ها را بررسی و ترکیب کند
		\item متن واحد را تدوین کرده و به \textbf{تصویب \lr{$\frac{2}{3}$} مجلس} برساند
		\item متن را به \textbf{رفراندوم عمومی} (۵۰\%+۱، حداقل مشارکت ۴۰\%) بگذارد
	\end{enumerate}
\end{bandbox}

\subsection{ب) شورای قضایی انتقالی}

\begin{bandbox}{بند ۱۵}
	مجلس مهستان از میان اساتید حقوق، اعضای کانون وکلا و قاضیان مستقل، \textbf{شورایی}
	برای رهبری دستگاه قضاییه تشکیل می‌دهد:
	\begin{itemize}
		\item تصفیه و پاکسازی دستگاه قضایی از عناصر فاسد
		\item بازبینی و تعلیق قوانین ناعادلانه
		\item سازمان‌دهی مجدد دادگاه‌ها
		\item تشکیل دادگاه‌های ویژه‌ی عدالت انتقالی
	\end{itemize}
\end{bandbox}

\subsection{ج) فرماندهی نظامی‌انتظامی موقت}

\begin{bandbox}{بند ۱۶}
	مجلس مهستان شورایی برای مدیریت نیروهای نظامی و انتظامی تعیین می‌کند.
	\textbf{فرمانده و جانشین فرمانده‌ی ستاد مشترک} مستقیماً توسط مجلس تعیین می‌شوند.
\end{bandbox}

\subsection{د) کمیسیون ملی اصلاح بخش امنیتی (\lr{SSR})}

\begin{bandbox}{بند ۱۷}
	کمیسیون تخصصی مستقل برای:
	\begin{itemize}
		\item غربال‌گری (\lr{Vetting}) فرماندهان و نیروهای امنیتی
		\item انحلال نهادهای سرکوب (سازمان عقیدتی‌سیاسی، گشت ارشاد)
		\item ادغام نیروهای واجد شرایط در ارتش واحد ملی
		\item واگذاری دارایی‌های اقتصادی نهادهای نظامی به صندوق ملی بازسازی
		\item آموزش حقوق بشر و تبعیت از مقام مدنی
	\end{itemize}
\end{bandbox}

\subsection{ه) کمیسیون عدالت انتقالی}

\begin{bandbox}{بند ۱۸}
	مجلس مهستان ظرف حداکثر \textbf{۳ ماه} طرح جامع عدالت انتقالی را تدوین خواهد کرد.
	در هر صورت، عدالت انتقالی شامل این عناصر خواهد بود:
\end{bandbox}
\vspace{0.5\baselineskip}
\noindent
\begin{center}
\begin{tikzpicture}[
    every node/.style={align=right, font=\small},
    node distance=0.5cm and 1.0cm,
]

% ----- CENTER NODE -----
\node[
    rectangle, rounded corners=8pt,
    draw=MahestanPurple, fill=MahestanLightPurple,
    text width=4cm, inner sep=8pt,
    font=\bfseries\small,
    align=right
] (center) {%
    {\centering \rl{عدالت انتقالی ایران}\par}
    \vspace{3pt}
    \rl{\footnotesize مدل نهایی را مجلس مهستان}\\
    \rl{\footnotesize با مشارکت مردم تعیین می‌کند}%
};

% ----- TOP NODE -----
\node[
    mahbox=MahestanBlue,
    yshift=45mm
] (tj1) {%
    {\centering \textbf{\rl{حقیقت‌یابی}}\par}
    \vspace{3pt}
    \rl{کمیسیون حقیقت‌یابی}\\
    \rl{و آشتی ملی}\\
    \rl{مستندسازی جنایات}\\
    \rl{ثبت شهادت قربانیان}%
};

% ----- RIGHT-UP NODE -----
\node[
    mahbox=MahestanRed,
    right=0.7cm of center,
    yshift=17mm
] (tj2) {%
    {\centering \textbf{\rl{عدالت کیفری}}\par}
    \vspace{3pt}
    \rl{محاکمات عادلانه}\\
    \rl{تفکیک سطوح مسئولیت}\\
    \rl{عفو مشروط (سطوح پایین)}%
};

% ----- RIGHT-DOWN NODE -----
\node[
    mahbox=MahestanOrange,
    right=0.7cm of center,
    yshift=-2.2cm
] (tj3) {%
    {\centering \textbf{\rl{جبران خسارت}}\par}
    \vspace{3pt}
    \rl{غرامت مالی}\\
    \rl{اعاده‌ی حیثیت}\\
    \rl{خدمات روانشناختی}%
};

% ----- LEFT-UP NODE -----
\node[
    mahbox=MahestanGold,
    left=0.7cm of center,
    yshift=17mm
] (tj4) {%
    {\centering \textbf{\rl{آشتی ملی و حافظه}}\par}
    \vspace{3pt}
    \rl{گفتگوهای ملی آشتی}\\
    \rl{موزه‌ها و یادبودها}\\
    \rl{اصلاح تاریخ‌نگاری}%
};

% ----- LEFT-DOWN NODE -----
\node[
    mahbox=MahestanGreen,
    left=0.7cm of center,
    yshift=-2.2cm
] (tj5) {%
    {\centering \textbf{\rl{تضمین عدم تکرار}}\par}
    \vspace{3pt}
    \rl{اصلاح نهادی}\\
    \rl{آموزش حقوق بشر}\\
    \rl{نهادهای نظارتی مستقل}%
};

% ----- ARROWS -----
\draw[maharrow=MahestanBlue, shorten <=4pt, shorten >=4pt] (center) -- (tj1);
\draw[maharrow=MahestanRed, shorten <=4pt, shorten >=4pt] (center) -- (tj2);
\draw[maharrow=MahestanOrange, shorten <=4pt, shorten >=4pt] (center) -- (tj3);
\draw[maharrow=MahestanGold, shorten <=4pt, shorten >=4pt] (center) -- (tj4);
\draw[maharrow=MahestanGreen, shorten <=4pt, shorten >=4pt] (center) -- (tj5);

\end{tikzpicture}
\end{center}
\vspace{0.75\baselineskip}

\subsection{و) نهاد ناظر مستقل انتخابات}

\begin{bandbox}{بند ۱۹}
	هیأتی مستقل متشکل از حقوقدانان، نمایندگان جامعه‌ی مدنی و ناظران بین‌المللی،
	مسئول نظارت بر تمام رأی‌گیری‌ها و رفراندوم‌های دوران گذار خواهد بود.
\end{bandbox}

\subsection{ز) آمبودزمان (نهاد حقوق شهروندی)}

\begin{bandbox}{بند ۲۰}
	نهاد مستقل رسیدگی به شکایات شهروندان از نهادهای دولتی. ریاست آن توسط مجلس مهستان
	انتخاب شده و مستقل از کابینه عمل می‌کند.
\end{bandbox}

\subsection{ح) کمیسیون ملی رسانه‌ها}

\begin{bandbox}{بند ۲۱}
	نهاد مستقل برای تبدیل صداوسیمای دولتی به \textbf{رسانه‌ی عمومی مستقل} (الگوی
	\lr{BBC/ARD})، تضمین آزادی مطبوعات، جلوگیری از انحصار رسانه‌ای و رفع کامل
	فیلترینگ و سانسور.
\end{bandbox}

\subsection{ط) شورای اقتصادی دوران گذار}

\begin{bandbox}{بند ۲۲}
	شورای مستقل متشکل از اقتصاددانان برجسته‌ی داخل و خارج کشور برای مشاوره و طراحی
	سیاست‌های اقتصادی. گزارش‌های شورا هم‌زمان به مجلس و کابینه ارائه می‌شود.
\end{bandbox}

\subsection{ی) مرکز تحقیقات مجلس مهستان (بازوی تکنوکراتیک)}

\begin{bandbox}{بند ۲۳}
	نهاد کنونی مرکز تحقیقات مجلس با ریاست جدید، \textbf{بازوی تخصصی و تکنوکراتیک}
	مجلس مهستان خواهد بود:
	\begin{itemize}
		\item هماهنگی با متخصصان و پژوهشگران
		\item تهیه‌ی گزارش‌های فنی‌تخصصی (\lr{Fact Sheets}) برای هر تصمیم
		\item ارائه‌ی \textbf{گزینه‌های سیاستی} (نه توصیه‌ی یکجانبه) به نمایندگان
	\end{itemize}
	\begin{quotebox}
		\textbf{نکته‌ی کلیدی:} مرکز تحقیقات \textbf{بُعد واقعیت} (\lr{What is}) را تأمین
		می‌کند. تصمیم نهایی (\lr{What ought to be}) با نمایندگان منتخب و مردم است.
	\end{quotebox}
\end{bandbox}

% ──────────────────────────────────────────────────────────
\section{واحد ارتباط مردمی و دموکراسی مستقیم (بازوی دموکراتیک)}
% ──────────────────────────────────────────────────────────

\begin{bandbox}{بند ۲۴}
	به‌موازات بازوی تکنوکراتیک، \textbf{واحد مستقل ارتباط مردمی و دموکراسی مستقیم}
	تشکیل خواهد شد، زیر نظر مستقیم هیأت رئیسه‌ی مجلس.
\end{bandbox}

\vspace{0.5\baselineskip}
\noindent
\begin{center}
	\resizebox{0.97\linewidth}{!}{%
		\begin{tikzpicture}[
			every node/.style={align=right, font=\small},
			node distance=0.5cm and 0.8cm
			]
			% Main central node
			\node[
			rectangle, rounded corners=8pt,
			draw=MahestanBlue, fill=MahestanLightBlue!60,
			text width=4cm, inner sep=10pt,
			font=\bfseries\small
			] (center) {
				\rl{واحد ارتباط مردمی}\\
				\rl{و دموکراسی مستقیم}\\[4pt]
				\rl{\footnotesize بازوی دموکراتیک}\\
				\rl{\footnotesize زیر نظر هیأت رئیسه}
			};
			
			% Row 1
			\node[mahbox=MahestanBlue, above right=1.2cm and 1.5cm of center] (f1) {
				\textbf{\rl{۱. بستر دموکراسی مستقیم}}\\
				\rl{مدیریت پلتفرم دیجیتال}\\
				\rl{رأی‌گیری‌های موضوعی}\\
				\rl{رفراندوم‌های رسمی}\\
				\rl{نظرسنجی‌های مشورتی}
			};
			\node[mahbox=MahestanGreen, above left=1.2cm and 1.5cm of center] (f2) {
				\textbf{\rl{۲. اولویت‌های مردمی}}\\
				\rl{سامانه‌ی ثبت پیشنهاد}\\
				\rl{خطوط تلفنی رایگان}\\
				\rl{صندوق‌های فیزیکی}\\
				\rl{گزارش ماهانه: «صدای مردم»}
			};
			\node[mahbox=MahestanGold, right=1.5cm of center] (f3) {
				\textbf{\rl{۳. جلسات عمومی}}\\
				\rl{جلسات در شهر و روستا}\\
				\rl{پنل‌های شهروندی}\\
				\rl{گفتگوهای محله‌ای}
			};
			\node[mahbox=MahestanPurple, left=1.5cm of center] (f4) {
				\textbf{\rl{۴. نظرسنجی‌های علمی}}\\
				\rl{نظرسنجی‌های نمایندگی}\\
				\rl{سنجش اولویت‌های مردمی}\\
				\rl{انتشار عمومی نتایج}
			};
			\node[mahbox=MahestanOrange, below right=1.2cm and 1.5cm of center] (f5) {
				\textbf{\rl{۵. اطلاع‌رسانی و شفافیت}}\\
				\rl{پخش زنده‌ی جلسات}\\
				\rl{ترجمه به زبان‌های محلی}\\
				\rl{پادکست و ویدئوی توضیحی}
			};
			\node[mahbox=MahestanRed, below left=1.2cm and 1.5cm of center] (f6) {
				\textbf{\rl{۶ و ۷. مدیریت طومارها}}\\
				\rl{ثبت طومارهای مردمی}\\
				\rl{پیگیری آستانه‌ی تعیین‌شده}\\
				\rl{درخواست‌های عزل نماینده}
			};
			
			\draw[maharrow=MahestanBlue] (center) -- (f1);
			\draw[maharrow=MahestanGreen] (center) -- (f2);
			\draw[maharrow=MahestanGold] (center) -- (f3);
			\draw[maharrow=MahestanPurple] (center) -- (f4);
			\draw[maharrow=MahestanOrange] (center) -- (f5);
			\draw[maharrow=MahestanRed] (center) -- (f6);
		\end{tikzpicture}
	}%
\end{center}
\vspace{0.5\baselineskip}

% ──────────────────────────────────────────────────────────
\section{اقدامات اولیه‌ی فوری}
% ──────────────────────────────────────────────────────────

\begin{bandbox}{بند ۲۵ --- ۱۰ فرمان اضطراری}
	مجلس مهستان در اولین هفته‌ی تشکیل، ۱۰ فرمان اضطراری صادر می‌کند:
	
	\begin{enumerate}
		\item \textbf{آزادی فوری} تمامی زندانیان سیاسی
		\item \textbf{رفع کامل فیلترینگ} و سانسور اینترنتی
		\item \textbf{تعلیق قوانین حجاب اجباری} و سایر قوانین تبعیض‌آمیز علیه زنان
		\item \textbf{حق تجمع و تظاهرات مسالمت‌آمیز} بدون مجوز
		\item \textbf{برقراری آزادی مطبوعات} و لغو تمام محدودیت‌های رسانه‌ای
		\item \textbf{توقف اعدام‌ها} تا برپایی دادگاه‌های عادلانه
		\item \textbf{ممنوعیت شکنجه و بدرفتاری} با بازداشت‌شدگان
		\item \textbf{انحلال گشت ارشاد} و نهادهای سرکوب مشابه
		\item \textbf{آزادی فعالیت احزاب و تشکل‌های مدنی} بدون مجوز
		\item \textbf{لغو قوانین تبعیض‌آمیز} علیه اقلیت‌های قومی و مذهبی
	\end{enumerate}
\end{bandbox}

% ──────────────────────────────────────────────────────────
\section{اولویت‌بندی قانون‌گذاری}
% ──────────────────────────────────────────────────────────

\vspace{0.5\baselineskip}
\noindent
\begin{adjustbox}{max width=\linewidth}
	\begin{tabular}{@{}lp{12cm}@{}}
		\toprule
		\textbf{زمان‌بندی} & \textbf{اقدامات قانون‌گذاری} \\
		\midrule
		\rowcolor{DangerBg}
		\textbf{فوری --- هفته‌ی ۱-۲} &
		قانون آزادی مطبوعات؛ قانون احزاب؛ قانون تظاهرات؛ شفافیت مالی مقامات؛
		لغو قوانین تبعیض علیه زنان \\
		\addlinespace
		\rowcolor{WarningBg}
		\textbf{اولویت بالا --- هفته‌ی ۲-۴} &
		قانون بحران اقتصادی؛ حقوق اقلیت‌ها و زبان مادری؛ حفاظت از داده‌های شخصی \\
		\addlinespace
		\rowcolor{NodeBg}
		\textbf{اولویت متوسط --- هفته‌ی ۴-۸} &
		قانون عدالت انتقالی؛ کمیسیون حقیقت‌یابی؛ قانون \lr{SSR}؛
		استقلال قضایی؛ استقلال بانک مرکزی \\
		\addlinespace
		\rowcolor{SuccessBg}
		\textbf{اولویت عادی --- هفته‌ی ۸-۱۴} &
		محیط زیست؛ اصلاح آموزش؛ خصوصی‌سازی شفاف؛ مجلس مؤسسان؛
		الحاق به کنوانسیون‌های بین‌المللی \\
		\bottomrule
	\end{tabular}
\end{adjustbox}
\vspace{0.5\baselineskip}

% ──────────────────────────────────────────────────────────
\section{فرایند تقنینی --- ترکیب تخصص و دموکراسی مستقیم}
% ──────────────────────────────────────────────────────────

\begin{bandbox}{بند ۲۶}
	فرایند تصویب هر قانون در مجلس مهستان از الگوی زیر پیروی می‌کند:
\end{bandbox}

\vspace{0.5\baselineskip}
\begin{center}
	\resizebox{0.9\linewidth}{!}{%
		\begin{tikzpicture}[
			every node/.style={align=right, font=\small},
			node distance=0.7cm
			]
			\node[mahbox=MahestanBlue, text width=6cm] (s1) {
				\textbf{\rl{۱. طرح یا لایحه مطرح می‌شود}}\\
				\rl{طرح: از سوی نمایندگان (حداقل ۱۰ امضا)}\\
				\rl{لایحه: از سوی کابینه}\\
				\rl{درخواست مردمی: طومار با آستانه‌ی تعیین‌شده}
			};
			\node[mahbox=MahestanGray, text width=6cm, below=0.5cm of s1] (s2) {
				\textbf{\rl{۲. ارجاع به کمیسیون تخصصی}}\\
				\rl{+ درخواست همزمان از دو بازو:}
			};
			\node[mahbox=MahestanBlue, text width=5.5cm, below right=0.5cm and 2cm of s2] (s2a) {
				\textbf{\rl{بازوی تکنوکراتیک}}\\
				\rl{گزارش فنی: داده‌ها، تجربه‌ی جهانی،}\\
				\rl{پیامدهای اقتصادی/حقوقی/اجتماعی،}\\
				\rl{حداقل ۳ گزینه‌ی سیاستی}
			};
			\node[mahbox=MahestanGreen, text width=5.5cm, below left=0.5cm and 2cm of s2] (s2b) {
				\textbf{\rl{بازوی دموکراتیک}}\\
				\rl{گزارش مردمی: انتشار عمومی خلاصه،}\\
				\rl{بازخوردگیری ۷-۳۰ روزه، نظرسنجی،}\\
				\rl{«مردم چه می‌خواهند؟»}
			};
			\node[mahbox=MahestanPurple, text width=6cm, below=1.8cm of s2] (s3) {
				\textbf{\rl{۳. بررسی در کمیسیون تخصصی}}\\
				\rl{با در اختیار داشتن هر دو گزارش}\\
				\rl{جلسه‌ی استماع + دعوت از متخصصان}
			};
			\node[mahbox=MahestanGray, text width=6cm, below=0.5cm of s3] (s4) {
				\textbf{\rl{۴. ارائه به صحن مجلس}}\\
				\rl{بحث آزاد + پیشنهاد اصلاحیه}
			};
			\node[mahbox=MahestanGold, text width=6cm, below=0.5cm of s4] (s5) {
				\textbf{\rl{۵. رأی‌گیری}}\\
				\rl{قوانین عادی: اکثریت ساده}\\
				\rl{قوانین مهم: \lr{$\frac{2}{3}$} آرا}\\
				\rl{اصلاح قانون اساسی موقت: \lr{$\frac{3}{4}$} آرا}
			};
			\node[decisionnode, below=0.5cm of s5] (dec) {\rl{نیاز به رفراندوم؟}};
			\node[mahbox=MahestanRed, text width=5cm, below right=0.5cm and 1.5cm of dec] (s7a) {
				\textbf{\rl{رفراندوم عمومی}}\\
				\rl{نظارت بین‌المللی}\\
				\rl{نیاز به ۵۰\%+۱ آرا}\\
				\rl{حداقل مشارکت: ۴۰\%}
			};
			\node[mahbox=MahestanGreen, text width=5cm, below left=0.5cm and 1.5cm of dec] (s7b) {
				\textbf{\rl{ابلاغ به کابینه برای اجرا}}
			};
			\node[mahbox=MahestanGreen, text width=6cm, below=2.2cm of dec] (s8) {
				\textbf{\rl{اجرا و نظارت}}\\
				\rl{نظارت مجلس + بازخورد مردمی}
			};
			
			\draw[maharrow=MahestanBlue] (s1) -- (s2);
			\draw[maharrow=MahestanBlue] (s2) -- (s2a);
			\draw[maharrow=MahestanGreen] (s2) -- (s2b);
			\draw[maharrow=MahestanBlue] (s2a) -- (s3);
			\draw[maharrow=MahestanGreen] (s2b) -- (s3);
			\draw[maharrow=MahestanGray] (s3) -- (s4);
			\draw[maharrow=MahestanGray] (s4) -- (s5);
			\draw[maharrow=MahestanGold] (s5) -- (dec);
			\draw[maharrow=MahestanRed] (dec) -- node[above]{\tiny\rl{بله}} (s7a);
			\draw[maharrow=MahestanGreen] (dec) -- node[above]{\tiny\rl{خیر}} (s7b);
			\draw[maharrow=MahestanGreen] (s7a) -- (s8);
			\draw[maharrow=MahestanGreen] (s7b) -- (s8);
		\end{tikzpicture}
	}%
\end{center}
\vspace{0.5\baselineskip}

% ──────────────────────────────────────────────────────────
\section{موارد الزامی رفراندوم}
% ──────────────────────────────────────────────────────────

\begin{bandbox}{بند ۲۷}
	موارد زیر \textbf{الزاماً} نیاز به رفراندوم عمومی خواهند داشت:
	\begin{itemize}
		\item تصویب قانون اساسی موقت
		\item تصویب قانون اساسی نهایی
		\item تعیین نوع حکومت (جمهوری/مشروطه)
		\item تعیین ساختار حکمرانی (متمرکز/فدرال/خودمختار)
		\item الحاق یا واگذاری سرزمینی
		\item هر موضوعی که مجلس مهستان با رأی \lr{$\frac{2}{3}$} آن را «سرنوشت‌ساز» تشخیص دهد
		\item هر موضوعی که به آستانه‌ی [X] امضای مردمی برسد
		\textit{(عدد دقیق توسط مجلس مهستان تعیین می‌شود)}
	\end{itemize}
\end{bandbox}

% ──────────────────────────────────────────────────────────
\section{سازوکارهای نظارت و پاسخ‌گویی}
% ──────────────────────────────────────────────────────────

\begin{bandbox}{بند ۳۰ --- عزل نماینده}
	در صورت عدم رضایت مردم یک منطقه از نماینده‌ی خود، با جمع‌آوری \textbf{۴۰\% از
		آرای مأخوذه در آن حوزه‌ی انتخاباتی در دور قبل} (نه ۴۰\% از واجدین شرایط)،
	نماینده عزل و فرد بعدی در فهرست جایگزین خواهد شد. نظارت بر جمع‌آوری امضا به
	عهده‌ی \textbf{هیأت مستقل نظارت بر انتخابات} است.
\end{bandbox}

\begin{bandbox}{بند ۳۱}
	در صورت فوت یا عدم توانایی مشارکت هر نماینده، فرد بعدی در فهرست از همان ناحیه
	جایگزین خواهد شد.
\end{bandbox}

\begin{bandbox}{بند ۳۲ --- مدت و تمدید مجلس مهستان}
	مجلس مهستان در اصل \textbf{نمی‌تواند بیش از دو سال} ادامه داشته باشد.
	\textbf{۲ ماه قبل از اتمام دو سال}، رأی‌گیری درباره‌ی تمدید یا خاتمه صورت می‌گیرد:
	
	\begin{itemize}
		\item اگر \textbf{فرایند گذار به پایان رسیده}: مجلس منحل و قدرت منتقل می‌شود.
		\item اگر \textbf{فرایند در جریان است}: مجلس می‌تواند با \textbf{رأی \lr{$\frac{2}{3}$} نمایندگان}،
		دوره‌ی خود را برای \textbf{حداکثر یک بار} و به مدت \textbf{حداکثر ۶ ماه} تمدید کند.
	\end{itemize}
	
	\begin{infobox}
		رأی‌گیری درباره‌ی تمدید دوره می‌تواند با \textbf{رفراندوم‌های دیگری که در همان
			مقطع برنامه‌ریزی شده} ترکیب شود تا هزینه‌های اجرایی کاهش یابد.
	\end{infobox}
\end{bandbox}

% ──────────────────────────────────────────────────────────
\section{مأموریت تدوین قانون اساسی نهایی}
% ──────────────────────────────────────────────────────────

\begin{bandbox}{بند ۳۳}
	مجلس مهستان پس از اجرای قانون اساسی موقت، پیش‌نویس \textbf{قانون اساسی نهایی}
	را از یکی از دو طریق تدوین می‌کند:
	\begin{itemize}
		\item \textbf{گزینه‌ی الف}: تشکیل \textbf{مجلس مؤسسان مستقل} (انتخابات جداگانه)
		\item \textbf{گزینه‌ی ب}: تفویض مأموریت به \textbf{کمیسیون ویژه‌ی درون مجلس}
		+ مشاوره‌ی عمومی گسترده
	\end{itemize}
	در هر دو حالت، قانون اساسی نهایی ملزم به رعایت \textbf{اصول غیرقابل‌نقض} مندرج
	در قانون اساسی موقت بوده و الزاماً از طریق \textbf{رفراندوم عمومی} تصویب خواهد شد.
\end{bandbox}

\begin{bandbox}{بند ۳۴ --- حل اختلاف بین نهادی}
	در صورت بروز اختلاف جدی بین نهادهای دوران گذار، \textbf{هیأت میانجی‌گری} متشکل
	از ۵ شخصیت مورد احترام ملی (پیشنهادی جامعه‌ی مدنی و مصوب مجلس) وارد عمل خواهد
	شد. در صورت عدم حل اختلاف، موضوع به \textbf{رفراندوم} گذاشته خواهد شد.
\end{bandbox}

% ============================================================
%  CHAPTER 4: INSTITUTIONAL STRUCTURE DIAGRAM
% ============================================================

% ── save portrait \textwidth for height constraint ──────────
\newlength{\portraittextwidth}
\setlength{\portraittextwidth}{\textwidth}

\begin{landscape}
\thispagestyle{plain}

% ── Chapter heading WITHOUT triggering \clearpage ───────────
% (landscape already issued \clearpage, so we place heading manually)
\refstepcounter{chapter}%
\addcontentsline{toc}{chapter}%
  {\protect\numberline{\thechapter}نمودار جامع ساختار نهادی}%
\label{ch:structure}%
\markboth{نمودار جامع ساختار نهادی}{نمودار جامع ساختار نهادی}%

{%
  \raggedleft
  \small\color{MahestanGray}%
  \chaptername\hspace{.5em}%
  \colorbox{MahestanDarkBlue}{\textcolor{white}{\thechapter}}%
  \par\nobreak\vspace{1ex}%
  {\color{MahestanBlue}\hrule height 0.8pt}\nobreak\vspace{0.4ex}%
  {\huge\bfseries\color{MahestanDarkBlue}%
   نمودار جامع ساختار نهادی}\par\nobreak%
  {\color{MahestanBlue}\hrule height 0.8pt}%
  \vspace{8pt}%
}%

% ── Diagram: constrained to fit remaining page ──────────────
% Available height ≈ landscape width − top/bottom margins − heading
% We cap both dimensions so it never overflows
\noindent
\begin{adjustbox}{max width=\linewidth, max height=\dimexpr\portraittextwidth-4cm\relax, center}
\begin{tikzpicture}[
  every node/.style={align=center, font=\footnotesize},
  node distance=0.6cm and 1cm,
  >={Stealth[length=4pt]}
]

%% ══════════════════════════════════════════════════════════
%%  ROW 0 — مردم ایران (top center)
%% ══════════════════════════════════════════════════════════
\node[
  rectangle, rounded corners=8pt,
  draw=MahestanGold, fill=MahestanGold!25,
  text width=5.5cm, inner sep=7pt,
  font=\small\bfseries
] (people) {
  \rl{مردم ایران}\\[2pt]
  {\footnotesize\rl{حاکمیت مطلق}}\\
  {\scriptsize\rl{انتخابات مستقیم --- رفراندوم‌ها --- نظرسنجی‌ها}}
};

% ── جامعه‌ی بین‌المللی (far left) ──
\node[
  rectangle, rounded corners=5pt,
  draw=MahestanGreen!40, fill=SoftGreen!50,
  text width=4cm, inner sep=4pt,
  left = 3cm of people
] (intl) {
  \textbf{\rl{جامعه‌ی بین‌المللی}}\\[2pt]
  {\scriptsize\rl{سازمان ملل --- اتحادیه اروپا}}\\
  {\scriptsize\rl{نظارت + حمایت فنی}}
};
% ── جامعه‌ی مدنی ──
\node[
  rectangle, rounded corners=5pt,
  draw=MahestanGreen!50, fill=SoftGreen,
  text width=4cm, inner sep=4pt,
  right=3cm of people
] (civil) {
  \textbf{\rl{جامعه‌ی مدنی}}\\[2pt]
  {\scriptsize\rl{نظارت + مشاوره}}
  {\scriptsize\rl{نقد + مطالبه}}
};
%% ══════════════════════════════════════════════════════════
%%  ROW 1 — بازوهای مجلس + مجلس مهستان (3-column)
%% ══════════════════════════════════════════════════════════

% ── مجلس مهستان (center) ──
\node[
  rectangle, rounded corners=8pt,
  draw=MahestanDarkBlue, fill=MahestanBlue!15,
  text width=5.5cm, inner sep=8pt,
  below=1.2cm of people,
  font=\small\bfseries
] (meh) {
  \rl{مجلس مهستان}\\[2pt]
  {\footnotesize\rl{نهاد عالی حاکمیت دوران گذار}}\\
  {\scriptsize\rl{۳۰۰ نماینده‌ی منتخب --- حداکثر ۲ سال}}
};

% ── واحد ارتباط مردمی (left of meh) ──
\node[
  rectangle, rounded corners=5pt,
  draw=MahestanBlue, fill=MahestanLightBlue!60,
  text width=3.2cm, inner sep=6pt,
  left=2cm of meh, yshift=1cm
] (demo) {
  \textbf{\rl{واحد ارتباط مردمی}}\\[2pt]
  {\scriptsize\rl{بازوی دموکراتیک}}\\
  {\scriptsize\rl{نظرسنجی --- بازخورد --- رفراندوم}}
};


% ── مرکز تحقیقات (right of meh) ──
\node[
  rectangle, rounded corners=5pt,
  draw=MahestanOrange, fill=MahestanLightOrange,
  text width=3.2cm, inner sep=6pt,
  right=2cm of meh, yshift=1cm
] (tech) {
  \textbf{\rl{مرکز تحقیقات}}\\[2pt]
  {\scriptsize\rl{بازوی تکنوکراتیک}}\\
  {\scriptsize\rl{واقعیت‌ها --- تحلیل فنی}}
};

%% ══════════════════════════════════════════════════════════
%%  ROW 2 — نهادهای مستقل + اجرایی (5-column spread)
%% ══════════════════════════════════════════════════════════

% ── هیأت نظارت بر انتخابات (far left) ──
\node[
  rectangle, rounded corners=5pt,
  draw=MahestanPurple!70, fill=MahestanLightPurple!50,
  text width=2.8cm, inner sep=5pt,
  below left=0.0cm and 4cm of meh
] (election) {
  \textbf{\rl{هیأت نظارت بر انتخابات}}\\[2pt]
  {\scriptsize\rl{مستقل}}
};
% ── آمبودزمان ──
\node[
  rectangle, rounded corners=5pt,
  draw=MahestanGold, fill=MahestanGold!10,
  text width=2.5cm, inner sep=4pt,
  below=1cm of election
] (ombud) {
  \textbf{\rl{آمبودزمان}}\\[2pt]
  {\scriptsize\rl{نهاد حقوق شهروندی}}
};

% ── شورای قضایی انتقالی (left) ──
\node[
  rectangle, rounded corners=5pt,
  draw=MahestanRed, fill=MahestanRed!10,
  text width=2.8cm, inner sep=5pt,
  below left=1.4cm and 0.5cm of meh
] (judicial) {
  \textbf{\rl{شورای قضایی انتقالی}}\\[2pt]
  {\scriptsize\rl{مستقل}}
};


% ── کابینه‌ی اجرایی موقت (right) ──
\node[
  rectangle, rounded corners=5pt,
  draw=MahestanGreen, fill=MahestanLightGreen,
  text width=2.8cm, inner sep=5pt,
  below right=1.4cm and 0.5cm of meh
] (cabinet) {
  \textbf{\rl{کابینه‌ی اجرایی موقت}}\\[2pt]
  {\scriptsize\rl{پاسخ‌گو به مجلس}}
};

% ── کمیسیون رسانه‌ها (far right) ──
\node[
  rectangle, rounded corners=5pt,
  draw=MahestanBlue!60, fill=NodeBg,
  text width=2.8cm, inner sep=5pt,
  below right=1.4cm and 4.5cm of meh
] (media) {
  \textbf{\rl{کمیسیون رسانه‌ها}}\\[2pt]
  {\scriptsize\rl{استقلال رسانه‌ای}}
};

%% ══════════════════════════════════════════════════════════
%%  ROW 3 — نهادهای زیرمجموعه (bottom row, 7-column)
%% ══════════════════════════════════════════════════════════

% ── کمیسیون حقیقت‌یابی ──
\node[
  rectangle, rounded corners=5pt,
  draw=MahestanPurple, fill=MahestanLightPurple,
  text width=2.5cm, inner sep=4pt,
  below=1.2cm of judicial
] (trc) {
  \textbf{\rl{کمیسیون حقیقت‌یابی}}\\[2pt]
  {\scriptsize\rl{عدالت انتقالی}}
};

% ── قانون اساسی موقت ──
\node[
  rectangle, rounded corners=5pt,
  draw=MahestanPurple, fill=MahestanLightPurple!40,
  text width=2.5cm, inner sep=4pt,
  below=1.5cm of meh
] (ctmp) {
  \textbf{\rl{قانون اساسی موقت}}\\[2pt]
  {\scriptsize\rl{ظرف ۲ ماه --- رفراندوم}}
};

% ── قانون اساسی نهایی ──
\node[
  rectangle, rounded corners=5pt,
  draw=MahestanPurple!80, fill=MahestanLightPurple!70,
  text width=3cm, inner sep=4pt,
  below=1.5cm of ctmp
] (cfin) {
  \textbf{\rl{قانون اساسی نهایی}}\\[2pt]
  {\scriptsize\rl{مجلس مؤسسان --- رفراندوم}}
};

% ── فرماندهی نظامی ──
\node[
  rectangle, rounded corners=5pt,
  draw=MahestanGray, fill=MahestanLightGray,
  text width=2.5cm, inner sep=4pt,
  below=1.2cm of cabinet
] (military) {
  \textbf{\rl{فرماندهی نظامی‌انتظامی}}\\[2pt]
  {\scriptsize\rl{موقت --- زیر نظر مجلس}}
};

% ── کمیسیون اصلاح امنیتی ──
\node[
  rectangle, rounded corners=5pt,
  draw=MahestanGray!80, fill=MahestanLightGray!60,
  text width=2.5cm, inner sep=4pt,
  right=0.4cm of military
] (ssr) {
  \textbf{\rl{کمیسیون اصلاح امنیتی}}\\[2pt]
  {\scriptsize\rl{\lr{SSR} --- غربال‌گری}}
};

% ── شورای اقتصادی (far right) ──
\node[
  rectangle, rounded corners=5pt,
  draw=MahestanGreen!70, fill=SuccessBg,
  text width=2.5cm, inner sep=4pt,
  below right = 0.5cm and 0.5cm of tech
] (econ) {
  \textbf{\rl{شورای اقتصادی گذار}}\\[2pt]
  {\scriptsize\rl{اقتصاددانان مستقل}}
};



%% ══════════════════════════════════════════════════════════
%%  ARROWS — all behind nodes via background layer
%% ══════════════════════════════════════════════════════════
\begin{scope}[on background layer]

  % ── Primary link: people → meh ──
  \draw[maharrow=MahestanGold, line width=1.8pt]
    (people) -- (meh)
    node[midway, right, xshift=3pt, font=\scriptsize]{\rl{انتخابات مستقیم}};

  % ── meh → arms ──
  \draw[maharrow=MahestanBlue]    (meh) -- (demo)
    node[midway, above, font=\tiny]{\rl{مدیریت}};
  \draw[maharrow=MahestanOrange]  (meh) -- (tech)
    node[midway, above, font=\tiny]{\rl{مدیریت}};

  % ── meh → row 2 institutions ──
  \draw[maharrow=MahestanPurple]  (meh) -- (election)
    node[midway, above left, font=\tiny]{\rl{تعیین}};
  \draw[maharrow=MahestanRed]     (meh) -- (judicial)
    node[midway, above left, font=\tiny]{\rl{تعیین}};
  \draw[maharrow=MahestanGreen]   (meh) -- (cabinet)
    node[midway, above right, font=\tiny]{\rl{تشکیل+نظارت}};
  \draw[maharrow=MahestanBlue]    (meh) -- (media)
    node[midway, above right, font=\tiny]{\rl{تعیین}};

  % ── meh → row 3 institutions ──
  \draw[maharrow=MahestanGold]    (meh) -- (ombud);
  \draw[maharrow=MahestanGray]    (meh) to[out=310, in=90] (military);
  \draw[maharrow=MahestanPurple]  (judicial) -- (trc);
  \draw[maharrow=MahestanGreen]   (cabinet) -- (military)
    node[midway, right, font=\tiny]{\rl{هماهنگی}};
  \draw[maharrow=MahestanGray]    (military) -- (ssr);

  % ── Constitution links (dashed) ──
  \draw[->, thick, MahestanPurple, dashed]
    (meh) to[out=250, in=90] (ctmp)
    node[near start, left, font=\tiny]{\rl{۲ ماه}};
  \draw[->, thick, MahestanPurple!80, dashed]
    (meh) to[out=280, in=90] (cfin);

  % ── Referendum links (dashed) ──
  \draw[->, thick, MahestanGold, dashed]
    (people) to[out=220, in=100] (ctmp)
    node[near end, left, font=\tiny]{\rl{رفراندوم}};
  \draw[->, thick, MahestanGold, dashed]
    (people) to[out=320, in=100] (cfin);

  % ── Feedback loops (dashed) ──
  \draw[->, thick, MahestanBlue, dashed]
    (demo) to[out=120, in=190] (people);
  \draw[->, thick, MahestanGreen!60, dashed]
    (civil) to[out=90, in=350] (meh);
  \draw[->, thick, MahestanGreen!50, dashed]
    (intl) to[out=60, in=210] (meh);
  \draw[->, thick, MahestanGreen!50, dashed]
    (econ) to[out=90, in=330] (meh);

\end{scope}

\end{tikzpicture}
\end{adjustbox}

\end{landscape}

% \chapter{نمودار جامع ساختار نهادی}
% \label{ch:structure}

% \begin{landscape}
% 	\thispagestyle{plain}
% 	\vspace*{0.3cm}
% 	\begin{center}
% 		{\large\bfseries\color{MahestanDarkBlue}\rl{نمودار جامع ساختار نهادی گذار}}\par
% 		\vspace{0.4cm}
% 		\resizebox{0.97\linewidth}{!}{%
% 			\begin{tikzpicture}[
% 				every node/.style={align=center, font=\small\rl},
% 				node distance=0.8cm and 1.2cm
% 				]
				
% 				%% ── مردم ایران (TOP) ──────────────────────────────────
% 				\node[
% 				rectangle, rounded corners=10pt,
% 				draw=MahestanGold, fill=MahestanGold!25,
% 				text width=6.5cm, inner sep=10pt, font=\bfseries
% 				] (people) {
% 					\rl{مردم ایران}\\[4pt]
% 					\rl{\footnotesize حاکمیت مطلق}\\
% 					\rl{\footnotesize انتخابات مستقیم --- رفراندوم‌ها --- نظرسنجی‌ها}
% 				};
				
% 				%% ── مجلس مهستان (CENTER) ──────────────────────────────
% 				\node[
% 				rectangle, rounded corners=10pt,
% 				draw=MahestanDarkBlue, fill=MahestanBlue!15,
% 				text width=6.5cm, inner sep=12pt,
% 				below=1.8cm of people, font=\bfseries\large
% 				] (meh) {
% 					\rl{مجلس مهستان}\\[4pt]
% 					\rl{\normalsize نهاد عالی حاکمیت دوران گذار}\\[2pt]
% 					\rl{\footnotesize ۳۰۰ نماینده‌ی منتخب --- حداکثر ۲ سال}
% 				};
				
% 				%% ── بازوی تکنوکراتیک (RIGHT TOP) ─────────────────────
% 				\node[
% 				rectangle, rounded corners=6pt,
% 				draw=MahestanOrange, fill=MahestanLightOrange,
% 				text width=4cm, inner sep=8pt,
% 				right=2.5cm of meh, yshift=1.5cm
% 				] (tech) {
% 					\textbf{\rl{مرکز تحقیقات}}\\[3pt]
% 					\rl{\footnotesize بازوی تکنوکراتیک}\\
% 					\rl{\scriptsize واقعیت‌ها --- تحلیل فنی}\\
% 					\rl{\scriptsize گزینه‌های سیاستی}
% 				};
				
% 				%% ── بازوی دموکراتیک (LEFT TOP) ───────────────────────
% 				\node[
% 				rectangle, rounded corners=6pt,
% 				draw=MahestanBlue, fill=MahestanLightBlue!60,
% 				text width=4cm, inner sep=8pt,
% 				left=2.5cm of meh, yshift=1.5cm
% 				] (demo) {
% 					\textbf{\rl{واحد ارتباط مردمی}}\\[3pt]
% 					\rl{\footnotesize بازوی دموکراتیک}\\
% 					\rl{\scriptsize نظرسنجی‌ها --- بازخورد}\\
% 					\rl{\scriptsize جلسات عمومی --- رفراندوم}
% 				};
				
% 				%% ── کابینه‌ی اجرایی موقت (RIGHT BOTTOM) ──────────────
% 				\node[
% 				rectangle, rounded corners=6pt,
% 				draw=MahestanGreen, fill=MahestanLightGreen,
% 				text width=3.8cm, inner sep=8pt,
% 				below right=1.8cm and 0.5cm of meh
% 				] (cabinet) {
% 					\textbf{\rl{کابینه‌ی اجرایی موقت}}\\[3pt]
% 					\rl{\footnotesize پاسخ‌گو به مجلس}
% 				};
				
% 				%% ── شورای قضایی انتقالی (LEFT BOTTOM) ────────────────
% 				\node[
% 				rectangle, rounded corners=6pt,
% 				draw=MahestanRed, fill=MahestanRed!10,
% 				text width=3.8cm, inner sep=8pt,
% 				below left=1.8cm and 0.5cm of meh
% 				] (judicial) {
% 					\textbf{\rl{شورای قضایی انتقالی}}\\[3pt]
% 					\rl{\footnotesize مستقل}
% 				};
				
% 				%% ── فرماندهی نظامی (BELOW CABINET) ───────────────────
% 				\node[
% 				rectangle, rounded corners=6pt,
% 				draw=MahestanGray, fill=MahestanLightGray,
% 				text width=3.5cm, inner sep=7pt,
% 				below=1.2cm of cabinet
% 				] (military) {
% 					\textbf{\rl{فرماندهی نظامی‌انتظامی}}\\[2pt]
% 					\rl{\footnotesize موقت --- زیر نظر مجلس}
% 				};
				
% 				%% ── کمیسیون اصلاح امنیتی (BELOW MILITARY) ────────────
% 				\node[
% 				rectangle, rounded corners=6pt,
% 				draw=MahestanGray!80, fill=MahestanLightGray!60,
% 				text width=3.5cm, inner sep=7pt,
% 				below=0.8cm of military
% 				] (ssr) {
% 					\textbf{\rl{کمیسیون اصلاح امنیتی}}\\[2pt]
% 					\rl{\footnotesize \lr{SSR} --- غربال‌گری نیروها}
% 				};
				
% 				%% ── کمیسیون حقیقت‌یابی (BELOW JUDICIAL) ──────────────
% 				\node[
% 				rectangle, rounded corners=6pt,
% 				draw=MahestanPurple, fill=MahestanLightPurple,
% 				text width=3.5cm, inner sep=7pt,
% 				below=1.2cm of judicial
% 				] (trc) {
% 					\textbf{\rl{کمیسیون حقیقت‌یابی و آشتی}}\\[2pt]
% 					\rl{\footnotesize عدالت انتقالی}
% 				};
				
% 				%% ── هیأت نظارت بر انتخابات (RIGHT OF DEMO) ───────────
% 				\node[
% 				rectangle, rounded corners=6pt,
% 				draw=MahestanPurple!70, fill=MahestanLightPurple!50,
% 				text width=3.5cm, inner sep=7pt,
% 				left=1.5cm of demo
% 				] (election) {
% 					\textbf{\rl{هیأت نظارت بر انتخابات}}\\[2pt]
% 					\rl{\footnotesize مستقل}
% 				};
				
% 				%% ── آمبودزمان (BELOW DEMO) ────────────────────────────
% 				\node[
% 				rectangle, rounded corners=6pt,
% 				draw=MahestanGold, fill=MahestanGold!10,
% 				text width=3.5cm, inner sep=7pt,
% 				below=3.2cm of demo
% 				] (ombud) {
% 					\textbf{\rl{آمبودزمان}}\\[2pt]
% 					\rl{\footnotesize نهاد حقوق شهروندی}
% 				};
				
% 				%% ── کمیسیون رسانه‌ها و شورای اقتصادی (RIGHT OF TECH) ─
% 				\node[
% 				rectangle, rounded corners=6pt,
% 				draw=MahestanBlue!60, fill=NodeBg,
% 				text width=3.5cm, inner sep=7pt,
% 				right=1.5cm of tech
% 				] (media) {
% 					\textbf{\rl{کمیسیون رسانه‌ها}}\\[2pt]
% 					\rl{\footnotesize استقلال رسانه‌ای}
% 				};
% 				\node[
% 				rectangle, rounded corners=6pt,
% 				draw=MahestanGreen!70, fill=SuccessBg,
% 				text width=3.5cm, inner sep=7pt,
% 				below=0.8cm of media
% 				] (econ) {
% 					\textbf{\rl{شورای اقتصادی گذار}}\\[2pt]
% 					\rl{\footnotesize اقتصاددانان مستقل}
% 				};
				
% 				%% ── قانون اساسی موقت و نهایی (BOTTOM CENTER) ─────────
% 				\node[
% 				rectangle, rounded corners=6pt,
% 				draw=MahestanPurple, fill=MahestanLightPurple!40,
% 				text width=3.5cm, inner sep=7pt,
% 				below=4cm of meh, xshift=-2.5cm
% 				] (ctmp) {
% 					\textbf{\rl{قانون اساسی موقت}}\\[2pt]
% 					\rl{\footnotesize ظرف ۲ ماه --- رفراندوم}
% 				};
% 				\node[
% 				rectangle, rounded corners=6pt,
% 				draw=MahestanPurple!80, fill=MahestanLightPurple!70,
% 				text width=3.5cm, inner sep=7pt,
% 				right=0.8cm of ctmp
% 				] (cfin) {
% 					\textbf{\rl{قانون اساسی نهایی}}\\[2pt]
% 					\rl{\footnotesize مجلس مؤسسان --- رفراندوم}
% 				};
				
% 				%% ── جامعه‌ی مدنی (FAR RIGHT BOTTOM) ──────────────────
% 				\node[
% 				rectangle, rounded corners=6pt,
% 				draw=MahestanGreen!50, fill=SoftGreen,
% 				text width=3cm, inner sep=6pt,
% 				below=3.8cm of tech, xshift=0.8cm
% 				] (civil) {
% 					\textbf{\rl{جامعه‌ی مدنی}}\\[2pt]
% 					\rl{\footnotesize نظارت + مشاوره}
% 				};
				
% 				%% ── جامعه‌ی بین‌المللی (FAR LEFT BOTTOM) ─────────────
% 				\node[
% 				rectangle, rounded corners=6pt,
% 				draw=MahestanGreen!40, fill=SoftGreen!50,
% 				text width=3cm, inner sep=6pt,
% 				below=1.5cm of election
% 				] (intl) {
% 					\textbf{\rl{جامعه‌ی بین‌المللی}}\\[2pt]
% 					\rl{\footnotesize سازمان ملل --- اتحادیه اروپا}\\
% 					\rl{\scriptsize نظارت + حمایت فنی}
% 				};
				
% 				%%─── Arrows ───────────────────────────────────────────────
% 				\draw[maharrow=MahestanGold, line width=2pt] (people) -- (meh)
% 				node[midway, right, xshift=4pt]{\footnotesize\rl{انتخابات مستقیم}};
% 				\draw[maharrow=MahestanOrange] (meh) -- (tech)
% 				node[midway, above]{\tiny\rl{مدیریت}};
% 				\draw[maharrow=MahestanBlue] (meh) -- (demo)
% 				node[midway, above]{\tiny\rl{مدیریت}};
% 				\draw[maharrow=MahestanGreen, thick] (meh) -- (cabinet)
% 				node[midway, above right]{\tiny\rl{تشکیل+نظارت}};
% 				\draw[maharrow=MahestanRed, thick] (meh) -- (judicial)
% 				node[midway, above left]{\tiny\rl{تعیین}};
% 				\draw[maharrow=MahestanGray] (meh) to[out=285, in=90] (military);
% 				\draw[maharrow=MahestanGray] (military) -- (ssr);
% 				\draw[maharrow=MahestanGreen] (cabinet) -- (military)
% 				node[midway, right]{\tiny\rl{هماهنگی}};
% 				\draw[maharrow=MahestanPurple] (meh) to[out=250, in=90] (trc);
% 				\draw[maharrow=MahestanPurple] (meh) to[out=200, in=0] (election);
% 				\draw[maharrow=MahestanGold] (meh) to[out=240, in=90] (ombud);
% 				\draw[maharrow=MahestanBlue] (meh) to[out=15, in=180] (media);
% 				\draw[maharrow=MahestanGreen] (meh) to[out=5, in=200] (econ);
% 				\draw[-Stealth, thick, MahestanPurple, dashed] (meh) to[out=260, in=90] (ctmp)
% 				node[near start, left]{\tiny\rl{۲ ماه}};
% 				\draw[-Stealth, thick, MahestanPurple!80, dashed] (meh) to[out=275, in=90] (cfin);
% 				\draw[-Stealth, thick, MahestanGold, dashed] (people) to[out=215, in=90] (ctmp)
% 				node[near end, left]{\tiny\rl{رفراندوم}};
% 				\draw[-Stealth, thick, MahestanGold, dashed] (people) to[out=200, in=110] (cfin);
% 				\draw[-Stealth, thick, MahestanBlue, dashed] (demo) to[out=240, in=90] (people);
% 				\draw[-Stealth, thick, MahestanGreen!60, dashed] (civil) to[out=90, in=270] (meh);
% 				\draw[-Stealth, thick, MahestanGreen!50, dashed] (intl) to[out=60, in=230] (meh);
% 			\end{tikzpicture}
% 		}%
% 	\end{center}
% \end{landscape}
% \vspace{0.5\baselineskip}

% ============================================================
%  CHAPTER 5: NON-NEGOTIABLE PRINCIPLES
% ============================================================
\chapter{اصول بنیادین غیرقابل‌نقض}
\label{ch:principles-fundamental}

این اصول خطوط قرمزی هستند که هیچ نهادی، حتی مجلس مهستان، حق تعلیق یا نقض آن‌ها را ندارد:

\begin{enumerate}
	\item \textbf{حاکمیت ملی}: قدرت از مردم و برای مردم
	\item \textbf{جدایی دین از دولت}: دولت در قبال تمام ادیان و مکاتب بی‌طرف است
	\item \textbf{برابری همه‌ی شهروندان}: صرف‌نظر از جنسیت، قومیت، مذهب و عقیده
	\item \textbf{آزادی‌های بنیادین}: آزادی بیان، مطبوعات، تجمع، تشکل، عقیده و مذهب
	\item \textbf{ممنوعیت مطلق شکنجه}: شکنجه و مجازات‌های خشن تحت هیچ شرایطی مجاز نیست
	\item \textbf{حق دادرسی عادلانه}: حق وکیل، محاکمه‌ی علنی، فرض بی‌گناهی و حق تجدیدنظر
	\item \textbf{حقوق اقلیت‌ها}: حقوق قومی، زبانی و مذهبی؛ حق آموزش به زبان مادری
	\item \textbf{حقوق زنان}: برابری کامل حقوقی، اجتماعی و اقتصادی
	\item \textbf{استقلال قضایی}: مستقل از قوای مجریه و مقننه
	\item \textbf{تمامیت ارضی}: ساختار حکمرانی از طریق رفراندوم تعیین می‌شود
	\item \textbf{حاکمیت قانون}: هیچ فرد یا نهادی بالاتر از قانون نیست
	\item \textbf{عدالت نه انتقام}: عدالت از طریق سازوکارهای قانونی
	\item \textbf{آزادی اطلاعات}: ممنوعیت سانسور و فیلترینگ
	\item \textbf{تعهد به صلح}: حل مسالمت‌آمیز اختلافات بین‌المللی
\end{enumerate}

\vspace{1em}
\subsection*{سازوکار تضمین اصول غیرقابل‌نقض}

\begin{infobox}
	در فاز اول گذار، ضمانت اجرایی این اصول عمدتاً بر \textbf{اراده و خواست مردم} استوار
	است که همین اصول را در رفراندوم قانون اساسی موقت تأیید کرده‌اند.
	
	پیشنهاد می‌شود \textbf{نهادهای مردم‌نهاد ناظر} برای پایش مستمر این اصول تشکیل شوند:
	\begin{itemize}
		\item گزارش‌دهی شفاف به مردم درباره‌ی هر مورد احتمالی نقض
		\item نقادی علنی و مستند از تصمیمات نهادهای گذار
		\item سنجش عملکرد مجلس مهستان و کابینه در قبال هر یک از ۱۴ اصل
	\end{itemize}
	
	در قانون اساسی نهایی، یک \textbf{دادگاه قانون اساسی مستقل} با صلاحیت ابطال قوانین
	مغایر با اصول بنیادین پیش‌بینی خواهد شد.
\end{infobox}

% ============================================================
%  CHAPTER 6: TIMELINE
% ============================================================
\chapter{خط زمانی جامع گذار}
\label{ch:timeline}

\begin{notebox}[title={\rl{یادداشت: زمان‌بندی نسبی}}]
	این سند در اسفند ۱۴۰۴ (فوریه‌ی ۲۰۲۶) تدوین شده است. تمام زمان‌بندی‌ها
	\textbf{نسبی} هستند و از \textbf{روز صفر} (آغاز رسمی گذار) محاسبه می‌شوند.
	تاریخ دقیق آغاز گذار از پیش قابل تعیین نیست.
\end{notebox}

\vspace{0.5\baselineskip}
\noindent
\begin{adjustbox}{max width=\linewidth}
	\begin{tabular}{@{}clp{8.5cm}@{}}
		\toprule
		\textbf{فاز} & \textbf{زمان‌بندی} & \textbf{اقدامات اصلی} \\
		\midrule
		\rowcolor{MahestanLightGray}
		\textbf{فاز ۰} & پیش از گذار &
		آمادگی، تدوین پیش‌نویس‌ها، ائتلاف‌سازی، زیرساخت فنی رأی‌گیری \\
		\rowcolor{WarningBg}
		\textbf{فاز ۱} & روز صفر تا هفته‌ی ۵ &
		انتخابات مجلس مهستان؛ تشکیل رسمی مجلس \\
		\rowcolor{DangerBg}
		\textbf{فاز ۲} & هفته‌ی ۱--۸ &
		۱۰ فرمان اضطراری؛ تشکیل کابینه و رأی اعتماد؛ مصوبات فوری \\
		\rowcolor{MahestanLightPurple!40}
		\textbf{فاز ۳} & هفته‌ی ۱--۱۰ &
		تدوین قانون اساسی موقت؛ رأی مجلس (\lr{$\frac{2}{3}$})؛ رفراندوم عمومی \\
		\rowcolor{NodeBg}
		\textbf{فاز ۴} & ماه ۳--۱۸ &
		نهادسازی؛ عدالت انتقالی؛ \lr{SSR}؛ اصلاحات اقتصادی \\
		\rowcolor{MahestanLightBlue!40}
		\textbf{فاز ۵} & ماه ۶--۱۸ &
		تدوین قانون اساسی نهایی؛ مشاوره‌ی گسترده؛ رفراندوم نهایی \\
		\rowcolor{SuccessBg}
		\textbf{فاز ۶} & ماه ۱۸--۲۴ (+۶) &
		انتخابات بر اساس قانون اساسی نهایی؛ انتقال رسمی قدرت \\
		\bottomrule
	\end{tabular}
\end{adjustbox}

\vspace{1cm}
\begin{center}
	\resizebox{\linewidth}{!}{%
		\begin{tikzpicture}[node distance=0.1cm]
			\def\bw{3.8cm}
			\def\bh{2.2cm}
			
			\node[phasenode=MahestanGray, text width=\bw, minimum height=\bh] (p0) {
				\textbf{\rl{فاز ۰}}\\[3pt]
				\rl{\footnotesize پیش از گذار}\\[2pt]
				\rl{\tiny آمادگی و ائتلاف‌سازی}
			};
			\node[phasenode=MahestanOrange, text width=\bw, minimum height=\bh, right=0.3cm of p0] (p1) {
				\textbf{\rl{فاز ۱}}\\[3pt]
				\rl{\footnotesize روز صفر + ۵ هفته}\\[2pt]
				\rl{\tiny انتخابات + تشکیل مجلس}
			};
			\node[phasenode=MahestanRed, text width=\bw, minimum height=\bh, right=0.3cm of p1] (p2) {
				\textbf{\rl{فاز ۲}}\\[3pt]
				\rl{\footnotesize هفته‌ی ۱--۸}\\[2pt]
				\rl{\tiny فرمان‌ها و اقدامات فوری}
			};
			\node[phasenode=MahestanPurple, text width=\bw, minimum height=\bh, right=0.3cm of p2] (p3) {
				\textbf{\rl{فاز ۳}}\\[3pt]
				\rl{\footnotesize هفته‌ی ۱--۱۰}\\[2pt]
				\rl{\tiny قانون اساسی موقت}
			};
			\node[phasenode=MahestanBlue, text width=\bw, minimum height=\bh, right=0.3cm of p3] (p4) {
				\textbf{\rl{فاز ۴}}\\[3pt]
				\rl{\footnotesize ماه ۳--۱۸}\\[2pt]
				\rl{\tiny نهادسازی و اصلاحات}
			};
			\node[phasenode=MahestanDarkBlue, text width=\bw, minimum height=\bh, right=0.3cm of p4] (p5) {
				\textbf{\rl{فاز ۵}}\\[3pt]
				\rl{\footnotesize ماه ۶--۱۸}\\[2pt]
				\rl{\tiny قانون اساسی نهایی}
			};
			\node[phasenode=MahestanGreen, text width=\bw, minimum height=\bh, right=0.3cm of p5] (p6) {
				\textbf{\rl{فاز ۶}}\\[3pt]
				\rl{\footnotesize ماه ۱۸--۲۴}\\[2pt]
				\rl{\tiny انتقال قدرت}
			};
			
			\foreach \a/\b in {p0/p1, p1/p2, p2/p3, p3/p4, p4/p5, p5/p6} {
				\draw[-Stealth, thick, MahestanGray] (\a) -- (\b);
			}
			
			% Timeline base line
			\draw[MahestanGold, line width=2pt]
			([yshift=-1.3cm]p0.south west) -- ([yshift=-1.3cm]p6.south east);
			
			% Milestone markers
			\foreach \phase/\mlabel in {
				p0/{آغاز آمادگی},
				p1/{آغاز گذار},
				p2/{تشکیل مجلس},
				p3/{قانون اساسی موقت},
				p6/{پایان گذار}%
			} {
				\node[circle, fill=MahestanGold, minimum size=8pt, inner sep=0]
				at ([yshift=-1.3cm]\phase.south) {};
				\node[below=0.15cm]
				at ([yshift=-1.3cm]\phase.south) {\tiny\rl{\mlabel}};
			}
		\end{tikzpicture}
	}%
\end{center}
\vspace{0.5\baselineskip}

% ============================================================
%  CHAPTER 7: CALL TO ACTION
% ============================================================
\chapter{فراخوان}
\label{ch:call}

از همه‌ی ایرانیان آزادیخواه دعوت می‌کنیم:

\begin{infobox}[title={\rl{به جریان‌ها، احزاب و ائتلاف‌ها}}]
	از هم‌اکنون پیش‌نویس قانون اساسی موقت پیشنهادی خود را تدوین و منتشر کنید.
\end{infobox}

\begin{infobox}[title={\rl{به هر شهروند}}]
	مطالعه کنید، بحث کنید، مشارکت کنید. نقد سازنده‌ی این طرح و هر طرح دیگری حق و وظیفه‌ی شماست.
\end{infobox}

\begin{infobox}[title={\rl{به متخصصان}}]
	دانش و تجربه‌ی خود را در خدمت طراحی نهادی گذار بگذارید.
\end{infobox}

\begin{infobox}[title={\rl{به فعالان حقوق بشر}}]
	مستندسازی را از هم‌اکنون ادامه دهید. این اسناد، مبنای عدالت انتقالی فردا خواهند بود.
\end{infobox}

\begin{summarybox}
	این بیانیه دعوتی است برای اندیشیدن، تکمیل و چکش‌کاری. سند زنده‌ای است که با
	مشارکت همگان بهتر خواهد شد. \textbf{تصمیم نهایی درباره‌ی آینده‌ی ایران فقط و فقط
		با مردم ایران است.}
\end{summarybox}